% ==================================================================
%  复习题 7
%  源：徐森林、薛春华《数学分析》第二册，印刷页 52-53 (PDF p. 58-59)
%  底本：PaddleOCR-VL-1.6 JSON/MD 识别结果，人工逐字核校与公式精修
%  规范：全西文标点，忠实原文条目编号与推导
% ==================================================================

\begin{exblock}{复习题 7}
  \item (杨忠道定理) 拓扑空间 $(X, \mathcal{T})$ 的任一子集 $A$ 的导集 $A'$ 为闭集 $\iff$ $(X, \mathcal{T})$ 的每个独点集的导集为闭集.

  \item 设 $(X, \mathcal{T})$ 为拓扑空间, $\{A_{\alpha} \mid \alpha \in \Gamma\}$ 为 $X$ 中的一个子集族, 如果 $\forall \alpha \in \Gamma$, 导集 $A_{\alpha}'$ 为 $(X, \mathcal{T})$ 中的闭集, 则集合 $\left(\bigcup\limits_{\alpha \in \Gamma} A_{\alpha}'\right)'$ 也为闭集.

  \item (1) 证明: $(\mathbb{R}^{n}, \mathcal{T}_{\rho_{0}}^{n})$ 中至多可数个稠密开集的交为 $(\mathbb{R}^{n}, \mathcal{T}_{\rho_{0}}^{n})$ 中的稠密开集;

  (2) 一般的度量空间 $(X, \mathcal{T}_{\rho})$, 上述结论仍正确吗?

  \item 设 $Y$, $Y_{\alpha}$ ($\alpha \in \Gamma$) 都是拓扑空间 $(X, \mathcal{T})$ 的连通子集. 如果 $\forall \alpha \in \Gamma$, $Y_{\alpha} \cap Y \neq \varnothing$, 则 $\left(\bigcup\limits_{\alpha \in \Gamma} Y_{\alpha}\right) \cup Y$ 为 $(X, \mathcal{T})$ 的连通子集. 上述「连通」改为「道路连通」结论如何? 并说明理由.

  \item 设 $\{Y_{\alpha} \mid \alpha \in \Gamma\}$ 为拓扑空间 $(X, \mathcal{T})$ 的道路连通的子集族, 并且满足: $\forall \alpha, \beta \in \Gamma$, 存在 $\Gamma$ 中的有限个元素 $\gamma_{1} = \alpha$, $\gamma_{2}$, $\cdots$, $\gamma_{n+1} = \beta$, 使得 $Y_{\gamma_{i}} \cap Y_{\gamma_{i+1}} \neq \varnothing$, $i = 1, 2, \cdots, n$, 则 $\bigcup\limits_{\alpha \in \Gamma} Y_{\alpha}$ 为道路连通子集. 上述「道路连通」改为「连通」结论如何? 并说明理由.

  \item 证明: 满足 $A_{1}$ 的可数紧致空间必为序列紧致空间.

  \item 举例说明拓扑空间中两个紧致子集的交可以不为紧致子集.

  \item 设 $(X, \rho)$ 为列紧度量空间, 则 $\forall \varepsilon > 0$, $X$ 必有 $\varepsilon$ 网 $A$ (即 $A$ 为 $X$ 的有限子集, 并且 $\forall x \in X$, 有 $\rho(x, A) = \inf\limits_{y \in A} \rho(y, x) < \varepsilon$).

  \item 设 $(X, \rho)$ 为度量空间, 证明:
  \[
    (X, \mathcal{T}_{\rho})\ \text{紧致} \iff (X, \mathcal{T}_{\rho})\ \text{列紧}.
  \]
  (提示: 应用 Lebesgue 数定理与本复习题第 8 题证明其充分性.)

  \item 设函数 $z = f(x, y)$ 在 $D = [0, 1] \times [0, 1]$ 上有定义, 且 $\forall x_{0} \in [0, 1]$, $f(x, y)$ 于 $(x_{0}, 0)$ 点连续, 证明: $\exists \delta > 0$, s. t. $f(x, y)$ 于 $D^{*} = \{(x, y) \mid 0 \leqslant x \leqslant 1,\ 0 \leqslant y \leqslant \delta\} = [0, 1] \times [0, \delta]$ 上有界.

  \item (1) 设 $u = f(x, y, z)$ 在闭立方体 $[a, b]^{3}$ 上连续, 证明: $g(x, y) = \max\limits_{a \leqslant z \leqslant b} f(x, y, z)$ 在正方形 $[a, b]^{2} \subset \mathbb{R}^{2}$ 上连续;

  (2) 设 $f(x, y, z)$ 在 $[a, b]^{3} = [a, b] \times [a, b] \times [a, b]$ 上连续, 令
  \[
    g(x) = \max_{a \leqslant y \leqslant x} \min_{a \leqslant z \leqslant b} f(x, y, z).
  \]
  证明: $g(x)$ 在 $[a, b]$ 上连续.

  \item 设 $U$ 为拓扑空间 $(X, \mathcal{T})$ 中的一个开集. $\mathcal{A}$ 为 $(X, \mathcal{T})$ 中的一个由紧致闭集构成的集族. 如果 $\mathcal{A}$ 满足 $\bigcap\limits_{A \in \mathcal{A}} A \subset U$, 则存在 $\mathcal{A}$ 的一个有限子族 $\{A_{1}, A_{2}, \cdots, A_{n}\}$ 也满足
  \[
    A_{1} \cap A_{2} \cap \cdots \cap A_{n} \subset U.
  \]

  \item 设 $(X, \mathcal{T})$ 为一个 Hausdorff 空间, $\mathcal{A}$ 为 $(X, \mathcal{T})$ 中由紧致子集构成的非空集族. 证明: 如果 $\mathcal{A}$ 中任意有限个元素的交是连通的, 则这个集族的交 $\bigcap\limits_{A \in \mathcal{A}} A$ 也是连通的.

  \item 设 $f(x, y)$ 满足: (1) 对固定的 $y \neq b$, $\lim\limits_{x \to a} f(x, y) = \psi(y)$; (2) $\exists \eta > 0$, s. t. $f(x, y)$ 当 $y \to b$ 时关于 $x \in E = \{x \mid 0 < |x - a| < \eta\}$ 存在一致极限 $\varphi(x)$. 证明:
  \[
    \lim_{x \to a} \lim_{y \to b} f(x, y) = \lim_{y \to b} \lim_{x \to a} f(x, y).
  \]

  \item 设 $f_{0}(x)$ 在 $[a, b]$ 上连续, $g(x, y)$ 在闭区域 $[a, b]^{2} = [a, b] \times [a, b]$ 上连续. 对任何 $x \in [a, b]$, 令
  \[
    f_{n}(x) = \int_{a}^{x} g(x, y)f_{n-1}(y)\,\mathrm{d}y, \quad n = 1, 2, \cdots.
  \]
  证明: 函数列 $\{f_{n}(x)\}$ 在 $[a, b]$ 上一致收敛于 $0$.

  \item 设有界点列 $\boldsymbol{z}^{n} = (x_{n}, y_{n})$, $n = 1, 2, \cdots$, 满足:
  \[
    \varliminf_{n \to +\infty} \|\boldsymbol{z}^{n}\| = l, \quad \varlimsup_{n \to +\infty} \|\boldsymbol{z}^{n}\| = L, \quad l < L,
  \]
  \[
    \lim_{n \to +\infty} \|\boldsymbol{z}^{n+1} - \boldsymbol{z}^{n}\| = 0.
  \]
  证明: 对 $\forall \mu \in (l, L)$, 圆周 $x^{2} + y^{2} = \mu^{2}$ 上至少有 $\{\boldsymbol{z}^{n}\}$ 的一个聚点.
\end{exblock}
